\section{Implementation and Experimental Protocol}
\label{sec:protocol}

\paragraph{Software.} The sampler is the Python package \texttt{kawasaki\_ising} (version 2.0.0), whose kernel is the one in \cref{alg:sweep}, compiled with Numba \citep{lam2015} on top of NumPy \citep{harris2020}. Experiments reported here were run on a two-core Intel Xeon @ 2.1\,GHz container with Python 3.13, NumPy 2.5.3, SciPy 1.18.1 and Numba 0.68.0. All experiment code is in \texttt{experiments/}, all raw outputs in \texttt{results/paper/}, and all tables and figures, together with most of the numbers quoted in the text, are generated from these files by \texttt{experiments/make\_paper\_assets.py}. The remaining numbers in the text are read off those tables.

\paragraph{Experiments.} The experiments are labelled E1--E6, E8 and E9. The check of the clustering metric (E8) is in \cref{app:metric}.
\begin{description}[leftmargin=1.2em,labelindent=0pt,itemsep=1pt]
\item[E1 (exhaustive local check).] For every one of the $\EoneDisp$ non-zero displacement vectors on the $6\times6$ torus and every assignment of spins to the union of the closed neighbourhoods of the two sites ($\le10$ sites), compare \eqref{eq:delta} and \eqref{eq:pair} with the difference of total energies before and after the exchange.
\item[E2 (exact stationary laws).] For the tori $3\times3$ ($M=3,4$) and $4\times4$ ($M=4,6$) enumerate $\Omega_{N,M}$, build $P_\beta$ and $\tilde P_\beta$ exactly (both for the production random-pair proposal and for nearest-neighbour proposals), solve $\mu P=\mu$ by sparse LU factorisation, and report $d_{\TV}(\mu,\pi_{\beta,M})$, the largest detailed-balance residual and the shift of the mean energy, at $T\in\{1,2,2.269,3\}$.
\item[E3 (implementation checks).] (a) Draw-for-draw equality of the experiment kernels with the package kernel; (b) the empirical rate of adjacent proposals against $4/N^2$; (c) reproducibility of seeded runs; (d) speed of the compiled kernel.
\item[E4 (impact on equilibrium observables).] At $f=1/2$ and $T\in\{1.5,\Tc,3\}$ run $R$ independent chains per rule ($R=400,200,100,40,20$ for $N=8,16,32,64,100$), $\EfiveBurn$ burn-in and $\EfiveSample$ sampling sweeps, and compare the replicate means of the energy per site $e=H/N^2$ between the two rules with a Welch standard error.
\item[E5 (variance peak versus exact binodal).] For $N\in\{32,64,100\}$, $f\in\{0.05,\allowbreak 0.10,\allowbreak 0.20,\allowbreak 0.30,\allowbreak 0.40,\allowbreak 0.50,\allowbreak 0.60,\allowbreak 0.80,\allowbreak 0.95\}$ and $\EfiveNT$ temperatures $1.5,1.55,\dots,3.0$, run $\EfiveR$ independent replicates of the estimator \eqref{eq:estimator} with $B=\EfiveBurn$, $S=\EfiveSample$ (the original protocol) and analyse them as described in \cref{sec:thermo}.
\item[E6 (single-seed protocol).] Re-run the protocol of the accompanying software for the variance-peak estimate (one seed per composition, $N=100$, a $20$-point grid on $[1.5,3.0]$) and compare with the exact binodal.
\item[E9 (equilibration).] E9a: long single traces at selected cells to compare the protocol window with late windows. E9b: a long equilibrium reference of the variance curves for six $(N,f)$ systems (\cref{sec:results-equil}).
\end{description}

\paragraph{Randomness.} Each run draws its seed from \texttt{numpy.random.SeedSequence} applied to the tuple identifying its (system size, composition, temperature, replicate, experiment) coordinates, so replicates are independent and every number is reproducible; a seed initialises both the NumPy generator used for the initial configuration and, via \texttt{seed\_all}, the generator used inside the compiled kernel.

\paragraph{Scope.} The proposal rule is the \emph{non-local} uniform random pair of the accompanying software. Its stationary law is the Gibbs law \eqref{eq:gibbs} (\cref{thm:correct}), but its dynamics are not physical migration dynamics, and we make no claim about relaxation times or coarsening kinetics.
